\documentclass[11pt,a4paper]{article}
\usepackage{amsmath,amssymb,amsthm}
\usepackage{geometry}
\usepackage{hyperref}
\usepackage{booktabs}
\geometry{margin=2.5cm}

\newtheorem{theorem}{Theorem}[section]
\newtheorem{corollary}[theorem]{Corollary}
\newtheorem{definition}[theorem]{Definition}
\newtheorem{remark}[theorem]{Remark}

\title{The Cascade Series --- Part I\\[0.5em]
\large The Cosmological Constant from the Observer's Frame}
\author{RTAC}
\date{March 2026}

\begin{document}
\maketitle

\begin{abstract}
Part 0~\cite{part0} derives the cascade invariant
$I_0 = \Omega_{19}\times\Omega_{217} = 1.2051\times 10^{-120}$ from pure
mathematics. This paper adds one input: the hypothesis that the cascade's
geometry is indistinguishable from our universe, and applies two
observer-frame corrections to convert the pure-number cascade invariant
into the dimensionless cosmological constant measured in physical
(reduced-Planck) units.

We observe four spacetime dimensions. The observer at $d = 4$ lives on
the $S^3$ boundary of $d = 5$, the volume maximum in Part 0's tower.
Two corrections connect the cascade invariant to the physical ratio:
(i)~the frame correction $(\Omega_5/\Omega_7)^2 = 9/\pi^2$ from the
cascade's natural reference at $d_0 = 7$ to the observer's host at
$d_V = 5$, and (ii)~the observer's cube--sphere bridge
$\Omega_2/V_3^{\mathrm{cube}} = 4\pi/8 = \pi/2$ at the spatial dimension
$d = 3$, which converts cascade-native sphere-area normalisation to the
cube-volume normalisation of the reduced Planck mass. Combined:
\[
\frac{\rho_\Lambda}{M^4_{\mathrm{Pl,red}}}
= \frac{2}{\pi}\cdot\frac{9\;\Omega_{19}\;\Omega_{217}}{\pi^2}
= \frac{18\;\Omega_{19}\;\Omega_{217}}{\pi^3}
= 6.996\times 10^{-121}.
\]
Observed (Planck 2018): $(7.150\pm 0.13)\times 10^{-121}$.  Leading
deviation: $-2.2\%$, placing $\Lambda$ in the descent-dependent
population. After the Part~0 Supplement first-order Gram correction
$\exp(\sum(1-C^2_{d,d+1}))=\exp(0.02108)$, the residual is $-0.07\%$,
well inside the $1.9\%$ Planck $1\sigma$ uncertainty.
\end{abstract}

\section{The Hypothesis}

\begin{definition}[Hypothesis]\label{def:hyp}
The infinite-dimensional unit ball, descended to four dimensions, is
indistinguishable from our universe.
\end{definition}

This is the single input beyond Part 0. It is tested by every prediction
in the series. It is not an identification, not a model, not an
approximation. It is a factual claim: no measurement can distinguish the
cascade's geometry from the observed universe.

\section{The Observer's Frame}

We observe four spacetime dimensions. The observer at $d = 4$ lives on the
$S^3$ boundary of $d = 5$. By Part 0's boundary dominance
($\Omega_{d-1}/V_d = d$), this boundary carries $4/5$ of the $d = 5$
content. The dimension $d_V = 5$ is the volume maximum of Part 0's tower:
the geometrically richest interior in the cascade.

Part 0's threshold pair ($\Omega_{19}$, $\Omega_{217}$) is generated by
the natural zero at $d_0 = 7$. The invariant
$I_0 = \Omega_{19}\times\Omega_{217}$ is intrinsically referenced to this
area maximum. The observer is not at $d_0 = 7$. The observer is at
$d_V = 5$.

Two frame corrections are required to connect the pure-number cascade
invariant to the physical cosmological-constant ratio measured in
reduced-Planck units.

\subsection{Correction 1: the host frame at $d_V = 5$}

The ratio of sphere areas at the two maxima:
\[
\frac{\Omega_5}{\Omega_7} = \frac{\pi^3}{\pi^4/3} = \frac{3}{\pi}.
\]
Applied once per threshold sphere area (both factors of the bilinear
invariant pass through the host rescaling; Part~0, Lemma~8.8), this
contributes $(\Omega_5/\Omega_7)^2 = 9/\pi^2$.

\subsection{Correction 2: the observer's cube--sphere bridge at $d = 3$}

Cascade content is sphere area (Part~0, Corollary~3.2: $\Omega_{d-1}$ is
the unique independent cascade quantity). The physical reduced Planck
mass $M_{\mathrm{Pl,red}}$ is a cube-volume normalisation: in the
zero-point mode counting that gives rise to $\rho_\Lambda$ in QFT, modes
are counted per unit $d^3 k/(2\pi)^3$ in momentum space, i.e., with
respect to a 3-cube measure. Cascade-native sphere-area normalisation
and $M_{\mathrm{Pl,red}}^4$ cube-volume normalisation differ by the
dimensionless ratio of sphere area to cube volume at the observer's
spatial dimension $d = 3$:
\[
\frac{\Omega_2}{V_3^{\mathrm{cube}}} \;=\; \frac{4\pi}{2^3} \;=\; \frac{\pi}{2}.
\]
Equivalently, converting the cube-normalised physical ratio into
cascade-native sphere units applies $\pi/2$; converting back applies
$2/\pi$.

\begin{remark}[Why $d = 3$ and not $d = 4$]
The cosmological constant is an energy density, not an energy; the
conversion factor between cascade sphere area and physical cube volume
lives at the observer's 3-dimensional spatial slice, not at the 4D
spacetime index. The observer's $S^3$ hosts the cascade's 3-sphere
surface $\Omega_2$ (the equatorial 2-sphere of the spatial slice); the
QFT cube is the 3-cube of momentum-space mode counting. Both are 3D
quantities at $d = 3$.
\end{remark}

\section{The Cosmological Constant}

\begin{theorem}\label{thm:cc}
\[
\frac{\rho_\Lambda}{M^4_{\mathrm{Pl,red}}}
\;=\; \underbrace{\frac{2}{\pi}}_{\substack{\text{observer bridge}\\d=3}}
\;\cdot\;\underbrace{\left(\frac{\Omega_5}{\Omega_7}\right)^{\!2}}_{\substack{\text{host frame}\\ d_V\to d_0}}
\;\cdot\;\Omega_{19}\;\Omega_{217}
\;=\; \frac{18\;\Omega_{19}\;\Omega_{217}}{\pi^3}
\;=\; 6.996\times 10^{-121}.
\]
\end{theorem}

\begin{proof}
The invariant $I_0 = \Omega_{19}\times\Omega_{217}$ is the product of
two threshold sphere areas (Part 0, Theorem~8.5). Each is referenced to
$d_0 = 7$. The observer at $d_V = 5$ re-references each through one
factor of $\Omega_5/\Omega_7 = 3/\pi$, contributing
$(3/\pi)^2 = 9/\pi^2$. Converting the cascade-native sphere-area
result to the reduced-Planck-mass cube-volume normalisation at the
observer's spatial dimension $d = 3$ multiplies by
$V_3^{\mathrm{cube}}/\Omega_2 = 2/\pi$ (equivalently, divides by the
sphere-to-cube ratio $\pi/2$):
\[
\frac{\rho_\Lambda}{M^4_{\mathrm{Pl,red}}}
= \frac{2}{\pi}\cdot\frac{9}{\pi^2}\times 1.20516\times 10^{-120}
= \frac{18}{\pi^3}\times 1.20516\times 10^{-120}
= 6.996\times 10^{-121}.\qedhere
\]
\end{proof}

Observed: $(7.150\pm 0.13)\times 10^{-121}$ (Planck 2018 $H_0=67.36$,
$\Omega_\Lambda=0.6847$; uncertainty $\pm 1.9\%$).

Leading deviation: $-2.2\%$. This places $\Lambda$ in the
\emph{descent-dependent} population of observables (negative sign at
leading order), alongside $\alpha_s$, $m_\tau/m_\mu$, $v$, $\ell_A$,
and $\Omega_m^{\mathrm{Bott}}$, all of which the Part~0 Supplement
first-order Gram correction closes from $\sim\!2\%$ to sub-$0.5\%$.
Applying the same correction to $\Lambda$:
\[
\frac{\rho_\Lambda}{M^4_{\mathrm{Pl,red}}}\bigg|_{\mathrm{Gram}}
= \frac{18}{\pi^3}\,\Omega_{19}\,\Omega_{217}\cdot\exp\!\bigl(\textstyle\sum_{d=5}^{216}(1-C^2_{d,d+1})\bigr)
= 6.996\times 10^{-121}\cdot e^{0.02108}
= 7.145\times 10^{-121}.
\]
Residual: $-0.07\%$, well inside the $\pm 1.9\%$ Planck $1\sigma$
uncertainty. The cosmological constant is not a Tier~1 structural
exception; it is a Tier~2 descent-dependent observable that closes to
sub-experimental precision through the same first-order correction
that closes the other descent-dependent members of the series.

Part 0 matches $\log_{10}|\Omega_{19}\Omega_{217}|$ to within $\sim 0.3$
without any frame correction. Both frame corrections ($9/\pi^2$ and
$2/\pi$) are Gamma-function identities at specified cascade dimensions,
not fits.

\section{The Formula Reads the Tower}

Part 0 establishes exactly four distinguished dimensions (Theorem~7.1).
The cosmological constant uses all four in the pure-cascade numerator
and adds one further factor --- the observer's cube--sphere bridge at
the spatial dimension $d = 3$ --- to convert to reduced-Planck units:
\[
\frac{\rho_\Lambda}{M^4_{\mathrm{Pl,red}}}
\;=\;
\underbrace{\frac{2}{\pi}}_{d=3\text{ bridge}}
\;\cdot\;
\underbrace{\frac{\Omega_{5}^{\,2}\;\Omega_{19}\;\Omega_{217}}{\Omega_7^{\,2}}}_{\text{Part 0 bilinear invariant}}
\;=\;
\frac{18\;\Omega_{19}\;\Omega_{217}}{\pi^3}.
\]

\begin{center}
\begin{tabular}{clll}
\toprule
Index & $d$ & Tower role & Formula role \\
\midrule
--- & 3 & Observer spatial slice & Cube--sphere bridge ($\pi/2$) \\
0 & 5 & Volume maximum & Observer's host (frame) \\
1 & 7 & Area maximum & Cascade reference (frame) \\
2 & 19 & Phase boundary & Near scale \\
3 & 217 & Terminus & Far scale \\
\bottomrule
\end{tabular}
\end{center}

Frame from the two maxima. Scale from the two thresholds. Bridge from
the observer's 3-dimensional spatial slice. Four distinguished cascade
dimensions plus the observer's spatial dimension.

\section{The Discrete Vacuum}

Between the two thresholds, the cascade has $217 - 19 = 198$ levels. Each
carries a specific sphere area, monotonically decreasing from
$\Omega_{18} = 0.59$ to $\Omega_{216}\approx 10^{-120}$. The vacuum is
198 discrete layers of geometry. The smallest fluctuation is
$\Omega_{217}\approx 10^{-120}$: the geometric floor, forced by the Gamma
function at $p(d) = \sqrt{\pi}$.

QFT integrates a continuous spectrum to the Planck cutoff and obtains
$\rho_\Lambda\sim M^4_{\mathrm{Pl}}$---$10^{120}$ too large. The cascade's
spectrum is discrete, finite, and terminates at $d_2 = 217$. The vacuum
energy is $10^{-120}$ because the vacuum has a floor.

\section{What This Paper Does and Does Not Do}

\textbf{Does:}
\begin{itemize}
\item Introduces the hypothesis (the single physical input of the series).
\item Derives
$\rho_\Lambda/M^4_{\mathrm{Pl,red}} = 18\Omega_{19}\Omega_{217}/\pi^3$
from Part~0's bilinear invariant together with two observer-frame
corrections: the host frame $(9/\pi^2)$ from $d_0=7$ to $d_V=5$, and the
cube--sphere bridge $(2/\pi)$ at the observer's spatial dimension
$d=3$. Leading deviation from Planck~2018: $-2.2\%$; after the Part~0
Supplement first-order Gram correction, residual $-0.07\%$, inside the
$\pm 1.9\%$ Planck $1\sigma$ uncertainty.
\item Places $\Lambda$ in the cascade's descent-dependent observable
population, closed by the same structural correction as $\alpha_s$,
$m_\tau/m_\mu$, $v$, $\ell_A$, and $\Omega_m^{\mathrm{Bott}}$.
\item Identifies the vacuum as 198 discrete layers with a natural terminus.
\item Identifies the cube--sphere bridge $\pi/2$ as the dimensional
conversion between cascade sphere-area content and QFT cube-volume mode
counting at the observer's 3D spatial slice.
\end{itemize}

\textbf{Does not:}
\begin{itemize}
\item Derive $d = 4$ from the cascade (used here as empirical fact).
\item Derive quantum mechanics, general relativity, the Standard Model, or
cosmological parameters beyond $\Lambda$.
\item Derive the Gram first-order correction from first principles; it
is imported from the Part~0 Supplement.
\end{itemize}

\textbf{Inputs:} One axiom (orthogonality, via Part 0). One hypothesis.
One empirical fact ($d = 4$). Zero free parameters.

\begin{thebibliography}{1}
\bibitem{part0} RTAC, \textit{The Cascade Series --- Part 0: Scale
Variance from Orthogonality}, 2026.
\end{thebibliography}

\end{document}